%% AMS-LaTeX Created with the Wolfram Language : www.wolfram.com

\documentclass[12pt]{article}
\usepackage[letterpaper,margin=1in]{geometry}
\usepackage{amsmath, amssymb, graphics, setspace}
\usepackage{graphicx} 
\graphicspath{{./figs/}}
\usepackage{tcolorbox}
\usepackage[framed,numbered,autolinebreaks,useliterate]{mcode}
\newcommand{\mathsym}[1]{{}}
\newcommand{\unicode}[1]{{}}

\usepackage[english]{babel}
\newtheorem{theorem}{Theorem}

\usepackage{caption}
\usepackage{subcaption}
\usepackage{bm}

\usepackage{tabulary}
\usepackage{multirow}
\usepackage{lscape}

\usepackage{hyperref}
\hypersetup{
    colorlinks=true,
    linkcolor=blue,
    filecolor=magenta,      
    urlcolor=cyan,
    pdftitle={Overleaf Example},
    pdfpagemode=FullScreen,
}

\def\mathbi#1{\textbf{\em #1}}

\begin{document}
\doublespacing	

\title{CE 401/5501: Applied Machine Learning\\ \large Topic 12-13: Unsupervised Learning - Lab Session}
\author{Dr. ZhiQiang Chen}
\date{}
\maketitle

\section*{Problem Statement}

Given: A moderate-size, high-dimensionality dataset \(\{(x, y)\}\), where \(x\) represents a data point with a high number of features (dimensions), denoted by \(d\), and \(y\) represents the corresponding label for the data point. The collection of all data points \(x\) is represented by a matrix \(X \in \mathbb{R}^{d \times N}\), where \(N\) is the total number of data points.

\subsection*{Objective}
\begin{itemize}
    \item \textbf{Data Exploration:} Gain an initial understanding of the data by visualizing relationships between features (e.g., using scatter plots of \(x_i\) vs. \(x_j\), where \(x_i\) indexes the \(i\)-th element of \(x\)). This exploration may reveal potential clustering trends within the data.
    \item \textbf{Clustering with K-means:} Implement the K-means clustering algorithm using a custom function that demonstrates the Expectation-Maximization (E-M) steps involved in iteratively assigning data points to clusters and updating cluster centroids. Observe the convergence of the summed cluster distances to the cluster centroids during the algorithm's execution.
    \item \textbf{Dimensionality Reduction and Visualization:} Apply dimensionality reduction techniques, specifically Principal Component Analysis (PCA) and t-distributed Stochastic Neighbor Embedding (t-SNE), to reduce the dimensionality of the data to 3. Visualize the reduced data to examine if clusters are apparent in the lower-dimensional space.
\end{itemize}

\subsection*{Formal Definition}

Given a dataset \(X \in \mathbb{R}^{d \times N}\) and corresponding labels \(y \in \{1, 2, \dots, K\}^N\), where \(d\) is the dimensionality of the data, \(N\) is the number of data points, and \(K\) is the number of clusters:
\begin{itemize}
    \item \textbf{Data Exploration:} Visualize pairwise relationships between features of \(X\) to identify potential clustering patterns.
    \item \textbf{Clustering with K-means:}
    \begin{itemize}
        \item Initialize \(k\) cluster centroids, \(c_1, c_2, \dots, c_k \in \mathbb{R}^d\).
        \item \textbf{E-step:} Assign each data point \(x_i\) to the cluster with the nearest centroid:
        \[
        \text{cluster}(x_i) = \arg\min_j \|x_i - c_j\|^2.
        \]
        \item \textbf{M-step:} Update each cluster centroid to the mean of the data points assigned to it:
        \[
        c_j = \frac{1}{|C_j|} \sum_{x_i \in C_j} x_i.
        \]
        \item Repeat the E-step and M-step until convergence (i.e., until the centroids no longer change significantly or a maximum number of iterations is reached).
    \end{itemize}
    \item \textbf{Dimensionality Reduction and Visualization:} Apply PCA and t-SNE to reduce the dimensionality of \(X\) to 3, resulting in \(X_{pca} \in \mathbb{R}^{3 \times N}\) and \(X_{tsne} \in \mathbb{R}^{3 \times N}\). Visualize \(X_{pca}\) and \(X_{tsne}\) to explore the presence of clusters in the lower-dimensional space.
\end{itemize}



\end{document}
